\subsection{Detailed proof of Theorem 3.1}
Fix for now $\alpha=w\cdot\alpha_i$. Set $K_\alpha:=K_{\alpha,1}$ and pick a sector-germ $\jepPubFrak{q}$ contained in $D_0^\alpha$. By (MA3), we know that for any $x\in\jepPubScript{I}$, there exists an apartment $A_x$ that contains both $x$ and $\jepPubFrak{q}$. Axiom (MA4) implies the existence of an isomorphism $\varphi_x:A_x\to\jepPubBbb{A}$ that fixes $\jepPubFrak{q}$, and~\cite[\S2.6]{jep88refRou11} ensures that $\varphi_x(x)$ does not depend on the choice of the apartment $A_x$ nor on the isomorphism $\varphi_x$, hence we can denote this element by $\rho_{\jepPubFrak{q}}(x)$. The map $\rho_q:\jepPubScript{I}\to\jepPubBbb{A}$ is the retraction of $\jepPubScript{I}$ onto $\jepPubBbb{A}$ centered at $\jepPubFrak{q}$, and its restriction to $\jepPubBbb{T}_\alpha$ does not depend on the choice of $\jepPubFrak{q}\in D_0^\alpha$.
\begin{remark}\label{jep88localRemark3_2}
Let $A$ be an apartment containing $\jepPubFrak{q}$ and $\varphi=\rho_q|_A$ be the restriction to $A$ of the retraction map $\rho_q$. Then $\varphi$ is the unique isomorphism of apartments that fixes $A\cap\jepPubBbb{A}$. Indeed, (MA4) implies the existence of an isomorphism of apartments $\psi:A\to\jepPubBbb{A}$ that fixes $\jepPubFrak{q}$. By definition, $\rho_{\jepPubFrak{q}}$ coincides with $\psi$ on $A$, hence $\varphi=\psi$ is an isomorphism of apartments, and fixes $\jepPubBbb{A}$, hence $\varphi$ fixes $\jepPubBbb{A}\cap A$. If $f:A\to\jepPubBbb{A}$ is an isomorphism of apartments that fixes $A\cap\jepPubBbb{A}$, then $f\circ\varphi^{-1}:\jepPubBbb{A}\to\jepPubBbb{A}$ is an isomorphism of affine spaces that fixes $\jepPubFrak{q}$, hence it must be trivial, which proves that $f=\varphi$ is unique.
\end{remark}
\begin{lemma}\label{jep88localLemma3_3}
Let $v\in\widetilde{K}_\alpha$ and $x\in\jepPubBbb{A}$ be such that $v\cdot x\in\jepPubBbb{A}$. Then $v$ fixes $D_{\lceil\alpha(x)\rceil}^\alpha$.
\end{lemma}
\begin{proof}
Set $A:=v\cdot\jepPubBbb{A}$ and let $k\in\jepPubBbb{Z}$ be such that $A\cap\jepPubBbb{A}$ contains $D_k^\alpha$. By~\cite[Lem.~3.2]{jep88refHeb16}, $A\cap\jepPubBbb{A}$ is a half-apartment,\footnote{Our definition of half-apartments is a bit different from the definition of~\cite{jep88refHeb16}: what we call half-apartments correspond to the true half-apartments of~\cite{jep88refHeb16}.} hence there exists an integer $k_1\in\jepPubBbb{Z}$ such that $A\cap A=D_{k_1}^\alpha$. Let $\varphi:\jepPubBbb{A}\to A$ be the isomorphism of apartments induced by $v$: Remark~\ref{jep88localRemark3_2} ensures that $\varphi$ fixes $A\cap\jepPubBbb{A}$, which means that $v$ fixes $A\cap\jepPubBbb{A}$. As $v\cdot x$ belongs to $A\cap\jepPubBbb{A}$, we obtain that $v\cdot(v\cdot x)=v\cdot x$, hence $v\cdot x=x$ belongs to $D_{k_1}^\alpha$. This implies that $D_{\lceil\alpha(x)\rceil}^\alpha$ is contained in $D_{k_1}^\alpha=A\cap\jepPubBbb{A}$, hence is fixed by $v$.
\end{proof}
\begin{lemma}\label{jep88localLemma3_4}
Let $x\in\jepPubBbb{A}$ and $M_x:=M_{\lceil\alpha(x)\rceil}^\alpha$. Then the map $f:K_\alpha\cdot x\to K_\alpha.M_x$ defined by $f(v\cdot x):=v\cdot M_x$ is well-defined and bijective.
\end{lemma}
\begin{proof}
Let $v,v'\in K_\alpha$ be such that $v\cdot x=v'\cdot x$: by Lemma~\ref{jep88localLemma3_3}, we get that $v^{-1}\cdot v'$ fixes $M_x$, hence $f$ is well-defined. Now assume that $v,v'\in K_\alpha$ satisfy $v\cdot M_x=v'.M_x$. Set $u:=v^{-1}v'$ and $A:=u\cdot\jepPubBbb{A}$: then $u\cdot M_x=M_x$ is contained in $A\cap\jepPubBbb{A}$ hence $u$ fixes $D_{\lceil\alpha(x)\rceil}^\alpha$ by Lemma~\ref{jep88localLemma3_3}. In particular, we have $u\cdot x=x$, i.e., $v\cdot x=v'\cdot x$ and $f$ is injective. As $f$ is surjective by definition, the lemma is proved.
\end{proof}
Set $\alpha_{\jepPubScript{I}}:=\alpha\circ\rho_{\jepPubFrak{q}}$ and, for any integer $k\geqslant0$, let $\jepPubScript{C}_k^\alpha$ be the set of all chambers $C$ that dominate some element of $K_\alpha.P_k$ and satisfy $\alpha_{\jepPubScript{I}}(C)>k$ (which means that there exists $X\in C$ such that $\alpha_{\jepPubScript{I}}(x)>k$ for all $x\in X$). Assume also that the chamber $C_k^\alpha$ chosen in Section~3.1 is not contained in $D_k^\alpha$.
\begin{lemma}\label{jep88localLemma3_5}
For any integer $k\geqslant0$, the map $g_k:K_\alpha M_{k+1}^\alpha\to\jepPubScript{C}_k^\alpha$ sending $v\cdot M_{k+1}^\alpha$ onto $v\cdot C_k^\alpha$ is well-defined and bijective.
\end{lemma}
\begin{proof}
The proof of the first part of the assertion is as in Lemma~\ref{jep88localLemma3_4}: if $v,v'\in K_\alpha$ are such that $v\cdot M_{k+1}^\alpha=v'.M_{k+1}^\alpha$, then $u:=v^{-1}\cdot v'$ satisfies $u\cdot M_{k+1}^\alpha\subset\jepPubBbb{A}$ and Lemma~\ref{jep88localLemma3_3} implies that $u\cdot C_k^\alpha=C_k^\alpha$, i.e., $v\cdot C_k^\alpha=v'\cdot C_k^\alpha$. As we moreover have $\alpha_{\jepPubScript{I}}(v\cdot C_k^\alpha)=\alpha(C_k^\alpha)>k$, $v\cdot C_k^\alpha=v'\cdot C_k^\alpha$ belongs to $\jepPubScript{C}_k^\alpha$ and the map $g_k$ is well-defined.

Assume now that $v,v'\in K_\alpha$ are such that $v\cdot C_k^\alpha=v'\cdot C_k^\alpha$. Set $u:=v^{-1}v'$ and let $X\in C_k^\alpha$ be an element fixed by $u$. Let $x\in X$ be such that $\alpha(x)>k$: by Lemma~\ref{jep88localLemma3_3}, $u$ fixes $M_{k+1}^\alpha\subset D_{\lceil\alpha(x)\rceil}^\alpha$, hence $g_k$ is injective.

It remains to check that $g_k$ is surjective, i.e., that $\jepPubScript{C}_k^\alpha=K_\alpha\cdot C_k^\alpha$. If $C\in\jepPubScript{C}_k^\alpha$, then there exists $u\in K_\alpha$ such that $C$ dominates $u\cdot P_k^\alpha$. By~\cite[Prop.~2.9.1)]{jep88refRou11}, there is an apartment $A$ that contains both $u\cdot D_k^\alpha$ and $C$, and Remark~\ref{jep88localRemark3_2} now gives an explicit isomorphism $\varphi:A\to\jepPubBbb{A}$ that fixes $A\cap\jepPubBbb{A}$. If $v\in K_\alpha$ induces $\varphi$, then $\alpha_{\jepPubScript{I}}(C)=\alpha(v\cdot C)$, hence $\alpha(v\cdot C)>k$, and $v\cdot C\subset\jepPubBbb{A}$ dominates $P_k^\alpha$, hence $v\cdot C=C_k^\alpha$, i.e., $C=v^{-1}\cdot C_k^\alpha$, which ends the proof.
\end{proof}
Combining Lemmas~\ref{jep88localLemma3_4} and~\ref{jep88localLemma3_5}, we get the following corollary.
\begin{corollary}\label{jep88localCorollary3_6}
For any $x\in\jepPubBbb{A}$, if $k:=\max(1,\lceil\alpha(x)\rceil)$, then $|K_\alpha\cdot x|=q_i'q_iq_i'\cdots$ ($k-1$ factors).
\end{corollary}
Until the end of this section, we assume that $W^v$ is infinite.
\begin{lemma}\label{jep88localLemma3_7}
Let $F$ be a type $0$ face of $\jepPubBbb{A}$. If $W^v$ is infinite, then there exists $g\in G$ such that $K_F/K_F\cap K_{g\cdot F}$ is infinite.
\end{lemma}
\begin{proof}
Let $g\in G$ be such that $a:=g\cdot0$ belongs to $C_f^v$ and set $F':=g\cdot F$. Let $(\alpha_n)_{n\geqslant0}$ be an injective sequence of positive real roots (i.e., $\alpha_n\in\Phi^+$ for any non-negative integer $n$). As we have $K_{\alpha_n}\subset K_F$, hence $|K_F\cdot F'|\geqslant|K_{\alpha_n}.a|$, for all $n\geqslant0$, it is enough to check (by Corollary~\ref{jep88localCorollary3_6} and thickness of $\jepPubScript{I}$) that $\alpha_n(a)\to+\infty$ as $n\to+\infty$.

By definition, any $\alpha_n$ can be written as $\sum_{i\in I}\lambda_{n,i}\alpha_i$ with $\lambda_{n,i}\in\jepPubBbb{Z}_+$ for all $(n,i)\in\jepPubBbb{Z}_+\times I$. The injectivity of the sequence $(\alpha_n)_{n\geqslant0}$ implies that $\sum_{i\in I}\lambda_{n,i}\to+\infty$ as $n$ goes to $+\infty$, hence $\lim_{n\to+\infty}\alpha_n(a)=+\infty$ as required.
\end{proof}
\begin{corollary}\label{jep88localCorollary3_8}
Let $F$ be a type $0$ face of $\jepPubScript{I}$. If $W^v$ is infinite, then there is no topology of topological group on $G$ for which $K_F$ is open and compact.
\end{corollary}
\begin{proof}
If there was such a topology, then for any $g\in G$, $K_F$ and $K_{g\cdot F}=gK_Fg^{-1}$ would be open and compact in $G$, hence $K_F\cap K_{g\cdot F}$ would have the same properties. This would imply the finiteness of the quotient $K_F/K_F\cap K_{g\cdot F}$ for any $g\in G$, which contradicts Lemma~\ref{jep88localLemma3_7}.
\end{proof}
